\begin{textsample}{NEG Dim 4 – global\_south\_2024\_05 – Score 1.00 – global\_south\_2024\_05/107933864.txt}  \label{ex:f4_neg_259}
India ' s Ministry of External Affairs ( MEA ) on Wednesday urged Indian students enrolled in US colleges and universities to abide by local laws and regulations with regard to the ongoing protests that have swept across American institutions for higher education against Israel ' s ongoing war in Gaza.
There have been no reports yet of the involvement of Indian students in these protests and the ministry has also said no student or their family has contacted Indian missions for help."We expect all our citizens at home and abroad to respect local laws and regulations,"MEA Spokesperson Randhir Jaiswal said at a news briefing in New Delhi in response to a question about protests at Columbia University.
Advertisement"So far no Indian students or their families have contacted us for any assistance in regard to disciplinary action which has been taken for their participation in the protest,"he added.
Hundreds of thousands of Indian students are enrolled in US colleges and universities in graduate and undergraduate courses.
Many of them desire to work after college and settle down here.
They @ @ @ @ @ @ @ @ @ @ an action that pits them against the law, which could jeopardize their future.
Protests have spread across US colleges and universities against Israel ' s ongoing war in Gaza, with students demanding an end to US support for Israel and, as in the case of Columbia University, cutting with businesses and other entities with investments in Israel.
These protests have turned violent in many instances with students taking over parts of the university building as they did in Columbia.
Police have arrested scores of students and are breaking up their encampments, as their tents pitched around the campuses are being called."Disclose, divest.
We will not slow, we will not rest"is what protestors at Columbia University have been heard chanting.
They want the university to divest in Israel and cut ties with companies that invest in Israel or have supported its war effort.
President Joe Biden addressed the violence in a speech from the White House saying that the freedom to protest must be accompanied by respect for the law @ @ @ @ @ @ @ @ @ @ is,"he said."It ' s against the law if violence occurs.
Vandalism, trespassing, breaking windows, shutting down campuses, forcing the cancellation of classes and graduations.
None of this is a peaceful protest.
Threatening people, intimidating people, instilling fear in people is not peaceful protest.
It ' s against the law.
Dissent is essential to democracy.
But dissent must never lead to disorder, or to denying the rights of others so students can finish the semester and their college education,"the US President said.
The US President also appealed for an end to anti-semitism and Islamophobia or discrimination against supporters of Palestine."There should be no place on any campus, no place in America, for anti-Semitism, or threats of violence against Jewish students.
There is no place for hate speech or violence of any kind, whether it ' s anti-semitism, Islamophobia, or discrimination against Arab-Americans or Palestinian Americans.
It ' s simply wrong.
There ' s no place for racism in @ @ @ @ @ @ @ @ @ @.
College authorities and law enforcement agencies have blamed some of the violence on outsiders joining student protestors.
New York police found"professional agitators"among protestors at Columbia.
Advertisement Related posts Revising its travel advisory for Iran and Israel, the government on Friday advised Indian nationals to remain vigilant while travelling to these countries and remain in touch with the Indian embassies there.

% matched lemmas: 
\end{textsample}
